%% =====================================================================
%%  Solutions du chapitre 1 — Hello, Python!
%% =====================================================================
\corrigesChapitre{1}{Hello, Python!}

\begin{corrigeatoi}{1}
\begin{console}
>>> 15 + 27
42
>>> 100 - 58
42
>>> 6 * 7
42
>>> print("Hello, Python!")
Hello, Python!
\end{console}
The three calculations give the same result, \py{42}\,---\,a little joke from the author. Fans of
the novel \emph{The Hitchhiker's Guide to the Galaxy} know that 42 is \enquote{the answer to the
ultimate question of life, the universe and everything}.
\end{corrigeatoi}

\begin{corrigeatoi}{2}
\begin{console}
>>> print("10 + 5")
10 + 5
>>> print(10 + 5)
15
>>> print("10", "+", 5)
10 + 5
>>> print(10, 5)
10 5
\end{console}
In the first instruction, \py{"10 + 5"} is between quotes: it is text, displayed as it is. In the
second, there are no quotes: Python calculates. The third displays three values separated by
spaces; it looks like the first one, but it is built differently. The last one displays two numbers,
separated by a space.
\end{corrigeatoi}

\begin{corrigeatoi}{3}
Here is an example; your own lines will of course be different.
\begin{code}{file=about_me.py}
print("Ana")
print("Lisbon")
print("Pasta with tomato sauce")
\end{code}
\begin{sortie}
Ana
Lisbon
Pasta with tomato sauce
\end{sortie}
\end{corrigeatoi}

\begin{corrigeatoi}{4}
The card is now 36 characters wide: two more dashes in the borders and two more spaces before each
right-hand bar. Every line still has the same length, so the frame stays straight.
\begin{code}{file=business_card.py}
# business_card.py - displays my business card
print("+----------------------------------+")
print("|                                  |")
print("|   SAM TAYLOR                     |")
print("|   Python learner                 |")
print("|                                  |")
print("|   sam@example.com                |")
print("|   Bristol, United Kingdom        |")
print("|                                  |")
print("+----------------------------------+")
\end{code}
\begin{sortie}
+----------------------------------+
|                                  |
|   SAM TAYLOR                     |
|   Python learner                 |
|                                  |
|   sam@example.com                |
|   Bristol, United Kingdom        |
|                                  |
+----------------------------------+
\end{sortie}
\end{corrigeatoi}

\begin{corrige}{1}
\demarche{look for quotes first: what is between quotes is displayed as it is; the rest is
calculated. Commas separate values, and \py{print()} puts one space between them.}
\begin{console}
>>> print("Python")
Python
>>> print("4 * 5")
4 * 5
>>> print(4 * 5)
20
>>> print("4", "*", 5)
4 * 5
>>> print(10 - 3, "days")
7 days
>>> print("Total:", 10 + 15)
Total: 25
\end{console}
\retiens{\py{print("4 * 5")} and \py{print("4", "*", 5)} display the same thing, but the first gives
one piece of text, the second three values separated by spaces.}
\end{corrige}

\begin{corrige}{2}
\demarche{Python runs the lines in order. \py{print()} with nothing inside displays an empty line,
and \py{3 + 2 + 1} is calculated before being displayed.}
\begin{code}{file=order.py}
print("3")
print("2")
print("1")
print()
print("Lift-off!", 3 + 2 + 1, "seconds later")
\end{code}
\begin{sortie}
3
2
1

Lift-off! 6 seconds later
\end{sortie}
\piege{forgetting the empty line: \py{print()} does display something\,---\,a line with nothing on
it.}
\end{corrige}

\begin{corrige}{3}
\begin{console}
>>> 125 + 389
514
>>> 1000 - 258
742
>>> 12 * 12
144
>>> 2 + 3 * 4
14
\end{console}
The cinema sold 514 tickets, 742 euros are left and the box holds 144 eggs.\par\vspace{1mm}
The last result is \py{14}, not \py{20}: Python did the multiplication \py{3 * 4} first, then added
\py{2}, exactly as in mathematics. Chapter~2 explains these priority rules in detail.
\end{corrige}

\begin{corrige}{4}
\demarche{try each wrong instruction in the shell and read the last line of each message.}
\begin{console}
>>> Print("Good morning")
Traceback (most recent call last):
  File "<pyshell#0>", line 1, in <module>
    Print("Good morning")
NameError: name 'Print' is not defined. Did you mean: 'print'?
>>> print("Good morning)
SyntaxError: unterminated string literal (detected at line 1)
>>> print(Good morning)
SyntaxError: invalid syntax. Perhaps you forgot a comma?
>>> print "Good morning"
SyntaxError: Missing parentheses in call to 'print'. Did you mean print(...)?
>>> pritn("Good morning")
Traceback (most recent call last):
  File "<pyshell#4>", line 1, in <module>
    pritn("Good morning")
NameError: name 'pritn' is not defined. Did you mean: 'print'?
>>> print("Good morning"))
SyntaxError: unmatched ')'
\end{console}
\begin{reponses}[2]
\item \py{NameError}: capital P.
\item \py{SyntaxError}: closing quote missing.
\item \py{SyntaxError}: quotes missing (see below).
\item \py{SyntaxError}: parentheses missing.
\item \py{NameError}: \py{pritn} is misspelt.
\item \py{SyntaxError}: one parenthesis too many.
\end{reponses}
The correct instruction is always the same: \py{print("Good morning")}.
\piege{in item~3 you may have expected a \py{NameError}, as with \py{print(Hello)}. But here there
are two words separated by a space, \py{Good} and \py{morning}, and Python cannot even make sense of
the sentence: it is a syntax error, detected before anything runs.}
\end{corrige}

\begin{corrige}{5}
\begin{code}{file=hello_me.py, cible=hello-me}
print("My name is Tom.")
print("I live in Manchester.")
print("My favourite food is pizza.")
\end{code}
\retiens{one \py{print()} per line displayed. Each piece of text starts and ends with the same kind of
quotation mark.}
\end{corrige}

\begin{corrige}{6}
\demarche{compare each line of the program with the line it must display. A piece of text needs
quotes; a number does not.}
\begin{code}{file=blanks.py, cible=blanks}
print("Good morning!")
print("Lunch is at", 12)
print("See you later!")
print(5, "+", 2, "=", 7)
\end{code}
\pourquoi{on line~4, \py{5} and \py{7} could also be written between quotes (\py{"5"} and
\py{"7"}): the display would be the same. But numbers are best written as numbers.}
\end{corrige}

\begin{corrige}{7}
\begin{reponses}
\item It is a \py{NameError}, on line~4: \py{prnit} instead of \py{print}. Python even suggests the
  right word.
\item The program is syntactically correct, so it started. Lines 1 to~3 ran normally. When Python
  reached line~4, it could not find \py{prnit}, stopped, and never ran line~5.
\item Replace \py{prnit} with \py{print}:
\end{reponses}
\begin{code}{file=menu.py}
print("Today's menu")
print("Starter: soup")
print("Main course: fish")
print("Dessert: apple pie")
print("Enjoy your meal!")
\end{code}
\begin{sortie}
Today's menu
Starter: soup
Main course: fish
Dessert: apple pie
Enjoy your meal!
\end{sortie}
\retiens{an exception stops the program at the faulty line; everything before it has already run.}
\end{corrige}

\begin{corrige}{8}
\demarche{run, read the message, fix one mistake, run again.}
\begin{etapes}
\item First run: \py{SyntaxError: unterminated string literal (detected at line 2)}. The closing
  quote is missing on line~2. Nothing is displayed, since it is a syntax error.
\item Second run: \py{SyntaxError: '(' was never closed}, on line~4: add the closing parenthesis.
\item Third run: the first two lines of the poem appear, then a \py{NameError} on line~3:
  \py{Print} has a capital letter.
\item Fourth run: the whole poem is displayed.
\end{etapes}
\begin{code}{file=poem.py, cible=poem}
print("Roses are red,")
print("Violets are blue,")
print("Python is fun,")
print("And so are you!")
\end{code}
\pourquoi{the two syntax errors were reported before anything ran, one at a time, in the order of
the lines. The \py{NameError} only appeared while the program was running.}
\end{corrige}

\begin{corrige}{9}
\demarche{copy each line of the drawing between quotes, including the spaces at the beginning:
they are what shifts the line to the right.}
\begin{code}{file=tree.py, cible=tree}
print("     *")
print("    ***")
print("   *o***")
print("  ***o***")
print(" *o*****o*")
print("    | |")
print("  =======")
\end{code}
\piege{forgetting the spaces at the beginning of the lines: the tree then leans to the left. Spaces
\emph{inside} the quotes are characters like any other.}
\end{corrige}

\begin{corrige}{10}
\begin{code}{file=ghosts.py}
# ghosts.py - what will appear?
print("one")
# print("two")
print("three")  # print("four")
#print("five")
print("six # seven")
\end{code}
\begin{sortie}
one
three
six # seven
\end{sortie}
\pourquoi{lines 1, 3 and~5 start with \py{\#}: they are comments. On line~4, everything after the
\py{\#} is a comment, so \py{"four"} is never displayed. On line~6, the \py{\#} is \emph{inside} the
quotes: it is part of the text, not the start of a comment.}
\end{corrige}

\begin{corrige}{11}
\demarche{each line displays a piece of text, then the result of a multiplication calculated by
Python.}
\begin{code}{file=table7.py, cible=table7}
print("7 x 1 =", 7 * 1)
print("7 x 2 =", 7 * 2)
print("7 x 3 =", 7 * 3)
print("7 x 4 =", 7 * 4)
print("7 x 5 =", 7 * 5)
\end{code}
\piege{using the letter \py{x} for the calculation: in Python the multiplication sign is \py{*}. The
\py{x} only appears inside the text.}
\end{corrige}

\begin{corrige}{12}
\demarche{the text of each line is fixed; the amounts are calculations. For the total, Python adds
the three products: \py{2 * 3 + 1 * 4 + 3 * 5}.}
\begin{code}{file=receipt.py, cible=receipt}
# receipt.py - Lucy's receipt at the Corner Cafe
print("====== CORNER CAFE ======")
print("2 x Coffee       =", 2 * 3)
print("1 x Muffin       =", 1 * 4)
print("3 x Orange juice =", 3 * 5)
print("-------------------------")
print("TOTAL =", 2 * 3 + 1 * 4 + 3 * 5, "euros")
print("Thank you, see you soon!")
\end{code}
\pourquoi{in the lines of the order, the spaces after \py{Coffee} and \py{Muffin} align the \py{=}
signs. \py{print()} adds a space between the text and the number, which is why there is no space
before the closing quote.}
\retiens{you can put calculations directly inside \py{print()}. In chapter~3, you will learn to
store the prices in variables, so that you only need to change them in one place.}
\end{corrige}

\begin{corrigequiz}
\begin{reponses}[5]
\item b
\item b
\item a
\item b
\item b
\item b
\item c
\item a
\item b
\item b
\end{reponses}
\pourquoi{4: the quotes make \py{3 + 4} a piece of text. 5: \py{print()} adds a space between the two
values. 9: a syntax error is detected before the first line runs. 10: the last line says what went
wrong; then the line number tells you where.}
\end{corrigequiz}
